\lab{6}{Оценка качества поиска}
\task
Оценить собственный поиск и две альтернативы на 30 запросах.
Вычислить P, DCG, NDCG и ERR на глубинах 1, 3 и 5; разобрать сильные
и слабые стороны систем и возможные улучшения.

\method
Собраны 30 запросов о памяти, синхронизации, процессах, файловых
дескрипторах и API. Набор включает короткие имена \code{fork},
\code{malloc}, \code{pthread_create}, словоформы
\code{memory allocation / memory allocations}, более широкие
\code{resource limits} и опечатку \code{epolll}. Он отражает несколько
типов сложности, но не является случайной выборкой пользовательского лога.

Сравниваются собственный строгий поиск, собственный TF-IDF и две
альтернативы: SQLite FTS5 и Tantivy. Готовые библиотеки используются
только для внешних базовых систем, собственный индекс ими не строится.
Выдачи сохраняются в TSV: система, запрос, ранг, URL. Релевантность
отдельно задаётся для пары «запрос, URL» по шкале 0--3: нерелевантный,
слабосвязанный, полезный, непосредственно отвечающий на запрос.
Для Precision релевантными считаются оценки не ниже 2.

Пусть \(g_i\) - оценка документа на позиции \(i\). Тогда:
\begin{align*}
P@k&=\frac1k\sum_{i=1}^{k}[g_i\ge2],\\
DCG@k&=\sum_{i=1}^{k}\frac{2^{g_i}-1}{\log_2(i+1)},\\
NDCG@k&=\frac{DCG@k}{IDCG@k},\\
ERR@k&=\sum_{i=1}^{k}\frac{R_i}{i}\prod_{j=1}^{i-1}(1-R_j),
\qquad R_i=\frac{2^{g_i}-1}{8}.
\end{align*}
Идеальный DCG получается сортировкой всех известных оценок запроса.
В оценочный набор входят разметка объединённого пула и дополнительные
эталонные документы. Поэтому IDCG не ограничен результатами одной системы.
При \(IDCG=0\) NDCG определён нулём. Короткая выдача дополняется
нулевыми оценками до глубины отсечения; её Precision всё равно делится
на \(k\). Неоценённый выданный документ принципиально отличается
от отсутствующего результата: программа завершает расчёт ошибкой.
Итоговые метрики усредняются по запросам, каждый запрос имеет одинаковый вес.

\journal
Сначала реализованы формулы и ручные контрольные примеры, затем
чтение TSV и проверка входных данных. Для сравнения построены индексы
SQLite и Tantivy на тех же 1\,679\,971 документах. Для каждого запроса
объединены первые пять результатов всех систем, повторные URL удалены.
Все 426 уникальных пар «запрос, документ» явно оценены с обоснованиями.
Оценивание учитывает не только заголовок, но и текст документа;
для неоднозначных длинных страниц дополнительно проверяются фрагменты
об искомом API. Контрольные суммы входных данных защищают от смешения
выдач и оценок разных экспериментов.

\testing
Ручной пример \(g=(3,0,2)\), идеальный список \((3,2)\), \(k=3\):
\[
P@3=\frac23,\quad DCG@3=8{,}5,\quad
NDCG@3=\frac{8{,}5}{7+3/\log_2 3},\quad
ERR@3=0{,}875+\frac{0{,}125\cdot0{,}375}{3}.
\]
Дополнительно проверяются пустая выдача, короткая выдача и NUL
в строке входа. Три модульных теста пройдены.
Дополнительные семь Python-тестов проверяют полноту пула,
отклонение пропущенных и повторных оценок, допустимость шкалы,
наличие обоснования, покрытие запросов и контрольные суммы.
Повторный расчёт метрик выполнен без чтения корпуса и индексов.

\results
\begin{table}[H]\centering
\begin{tabular}{lrrr}\toprule
Система & Результатов & Оценено & Покрытие, \%\\\midrule
Собственный булев поиск & 145 & 145 & 100,00\\
Собственный TF-IDF & 145 & 145 & 100,00\\
SQLite FTS5 & 145 & 145 & 100,00\\
Tantivy & 145 & 145 & 100,00\\\bottomrule
\end{tabular}
\caption{Покрытие оценками первых пяти результатов 30 запросов.}
\end{table}
145 вместо 150 результатов связано с пустой выдачей запроса
\code{epolll}. Он включён в усреднение и имеет нулевые метрики.
\begin{longtable}{llrrrr}
\caption{Средние метрики по 30 запросам.}\\
\toprule Система & Глубина & P & DCG & NDCG & ERR\\\midrule\endfirsthead
\toprule Система & Глубина & P & DCG & NDCG & ERR\\\midrule\endhead
Булев поиск & @1 & 0,1000 & 0,8333 & 0,1190 & 0,1042\\
Булев поиск & @3 & 0,1667 & 2,0484 & 0,1469 & 0,1942\\
Булев поиск & @5 & 0,1933 & 2,9820 & 0,1744 & 0,2283\\
\midrule
TF-IDF & @1 & 0,4333 & 2,9333 & 0,4190 & 0,3667\\
TF-IDF & @3 & 0,4333 & 5,2305 & 0,3557 & 0,4423\\
TF-IDF & @5 & 0,4933 & 7,8437 & 0,4219 & 0,4742\\
\midrule
SQLite FTS5 & @1 & 0,3667 & 2,1667 & 0,3095 & 0,2708\\
SQLite FTS5 & @3 & 0,4667 & 5,1214 & 0,3578 & 0,4257\\
SQLite FTS5 & @5 & 0,4067 & 6,5297 & 0,3745 & 0,4448\\
\midrule
Tantivy & @1 & 0,4000 & 2,4000 & 0,3429 & 0,3000\\
Tantivy & @3 & 0,4667 & 5,1563 & 0,3608 & 0,4421\\
Tantivy & @5 & 0,4267 & 6,9174 & 0,3937 & 0,4716\\
\bottomrule
\end{longtable}

На глубине 5 собственный TF-IDF имеет наибольшие значения всех
четырёх метрик среди сравниваемых систем. Его P@5 равна 0,4933
против 0,1933 строгого поиска, NDCG@5 равна 0,4219 против 0,1744.
При этом на глубине 3 Precision у SQLite FTS5 и Tantivy выше:
0,4667 против 0,4333. Поэтому преимущество зависит от глубины и метрики.

Разбор отдельных запросов показывает разные типы ошибок.
На \code{malloc} собственный TF-IDF достигает NDCG@5 = 1,0000,
а на \code{reference counting} получает 0,0000. В последнем случае
SQLite FTS5 даёт NDCG@5 = 0,4524. Общие слова и широкое объединение
кандидатов могут вытеснять документы, отвечающие именно на задачу
пользователя. Для исправления полезны учёт словосочетаний,
близости слов и специфики справочных страниц.

Оценки выставлены ассистентом без независимых человеческих оценщиков.
Кроме того, пул ограничен первыми пятью результатами сравниваемых
систем. Полное покрытие этого пула не означает полной разметки корпуса.

\conclusions
Выполнено сравнение четырёх систем на общем размеченном наборе.
TF-IDF существенно улучшил качество относительно порядка выдачи
строгого поиска, однако не оказался лучшим для каждого запроса
и каждой глубины. Чтобы повысить надёжность выводов, нужны
несколько оценщиков, расширение пула и более разнообразный набор запросов.
